\subsection{Functional Reactive Programming}\label{section:frp:frp_intro}

Reactive programs are interactive programs that communicate with the environment through inputs and outputs. Examples of such programs include servers, graphical user interfaces, and games.

Functional Reactive Programming (FRP) is a programming paradigm that combines functional programming with reactive programming. It allows developers to create programs that can react to changes in data over time, making it easier to handle asynchronous events and manage state in a declarative way. 

The main abstraction of FRP is the \textit{signal}, which models time-varying values. Conceptually, a signal is a sequence of values produced by interacting with the environment. Signals can be transformed, combined and manipulated using familiar functional techniques.

%Problems of FRP
However, representing signals as infinite streams makes it easy to write \todo{faulty} programs. Such programs may violate one of the three main properties \textit{causality} (computations are not dependent on the results of future computations), \textit{productivity} (computations eventually produce an output), and the \textit{absence of space leaks}. 

To overcome these issues, there exists a family of FRP languages~\cite{neel_krishnaswami_2013,bahr_asynchronous_2023,rizzo_modal_2025} which use \textit{modal types} to guarantee these properties. In literature~\cite{bahr_asynchronous_2023,rizzo_modal_2025}, these modal types are the `later' (written $\bigcirc a$) and the `box' (written $\square a$) type constructors. The type $\bigcirc a$ is a promise that a value of type $a$ will become available at some later time step, whereas the $\square a$ describes a value of type $a$ that \textit{time-independent} time steps and as result safe to keep across time steps.

Given these types it is possible to define a signal as a tuple of a current value (head) and a delayed computation (tail) for computing future values: $\text{type} \; \texttt{Signal }a = (a \times \texttt{Later }(\texttt{Signal }a))$.

\subsection{Rizzo}\label{section:frp:rizzo}

Rizzo~\cite{rizzo_modal_2025} is an FRP language that guarantees the same properties but removes the `box' type constructor by changing the semantics of signals. In particular, Rizzo redefines signals to be mutable references to a tuple of a head and tail:

\[\text{type} \texttt{ Signal }a = \text{Ref }(a \times \texttt{Later }(\texttt{Signal }a))\]

In Rizzo, signals are now updated in place whenever the environment produces a value that is relevant for the signal. These updates are done, not by the programmer, but instead by the reactive semantics of Rizzo, see Section~\ref{section:frp:rizzo}. For convenience we define the signal constructor as an infix operator: 
\[\_ :: \_ : a \rightarrow \texttt{Later }(\texttt{Signal }a) \rightarrow \texttt{Signal }a\]

Rizzo, like other languages~\cite{bahr_asynchronous_2023}, differentiates between two variants of the `later' type. 
In Rizzo there is the $\texttt{Later } a$ which describes a value of type $a$ becoming available at some (existential) future time step. 
And then there is the $\texttt{Delay }a$ which describe a value of type $a$ that is available at all (universally quantified) future time steps.

\todo{This delay type is how Rizzo can guarantee \textit{productivity}. It guards recursion with a delay, ensuring recursion cannot happen until a future time step is produced. This is useful in the formal model but a functional language without general recursion is inconvenient to work with and so going forward we assume general recursion. \dots I suspect this block is malplaced + also what are the ramifications?}

Rizzo has a number of constructors for working with delayed computations~\cite{rizzo_modal_2025}. We shall present some of them here but for a full overview \todo{see section?}. 
Starting with the delay type, its main constructor is $\text{delay} : a \rightarrow \texttt{Delay }a$ which simply pushes any value into the future. 
Then for the later type, we may use the constructor $\text{wait} : \text{Chan }a \rightarrow \texttt{Later }a$, which is used to wait for the environment to produce input on a particular channel. One can think of a channel ($\text{Chan}$ type) as an opening through which the environment can pass values to Rizzo programs.
Another later constructor is the $\text{tail} : \texttt{Signal }a \rightarrow \texttt{Later }(\texttt{Signal }a)$, which let us access the future value of a signal. 
However, probably the most useful of all is the applicative action on $\texttt{Later}$ computations:

\[f \triangleright x : \texttt{Delay }(a \rightarrow b) \rightarrow \texttt{Later }a \rightarrow \texttt{Later }b\]

This allows later values to be applied to delayed functions but the actual application does not happen until a future time step when the argument $x$ becomes available.

\todo{Figure with later modality types?}

With these constructors we already have the necessary tools to define some signal combinators of our own! 
The $mk\_sig$ function allows us to create signals from any delayed computation:

\[ \text{mk\_sig} \; later = (\lambda x. x :: \text{mk\_sig }later) \; \triangleright \; later \]

Assuming there exists a channel $console : \text{Chan String}$ that reads from the console, we can define signals that react to console input. An example of such a signal is $(\text{""} :: \text{mk\_sig} \; (\text{wait} \; console))$, which echoes the strings that are read from the console.

We can take it one step further and define the $\text{map} : (a \rightarrow b) \rightarrow \texttt{Signal }a \rightarrow \texttt{Signal }b$ function:

\[ \text{map} \; f \; (x :: xs) = f \; x :: (delay \; (map \; f)) \triangleright xs \]

The definition of map makes use of pattern matching on signals. Alternatively, we could have used explicit $\text{head}$ and $\text{tail}$ to extract $x$ and $xs$ respectively.

Let us return to what makes Rizzo different, its semantics of signals. Because signals are updated in place as mutable references, as opposed to stored immutable sequences is that it allows combinators like $\text{sample} : \texttt{Signal }a \rightarrow \texttt{Signal }b \rightarrow \texttt{Signal } (a \times b)$ defined as such:

\[\text{sample} \; xs \; ys = map \; (\lambda x. (x, head ys)) \; xs\]

The $sample$ combinator allows sample the current value of $ys$, every time that a new value $x$ is produced on the signal $xs$. If it weren't for the mutable semantics, this combinator would have had to keep the entirety of $ys$ (beginning from its first value) in memory.

Rizzo is an asynchronous FRP language which means different signals are not necessarily updated at the same time. This means input sources can produce values 

\subsection{Signals and the heap}\label{section:frp:signal_heap}

In Rizzo signals are stored in a heap structure, where the structure is a form of linked list. The importance of this structure is for the runtime to iterate and update signals for each input on a channel.

\todo{Skal man introducere lynhurtig signal-heap?}

Without going too much in detail, the way signals are stored in Rizzo is through a double linked list of signal nodes, where each node contains the current value of the signal and a pointer to the next node. When a signal is updated, a new node is created with the new value and the value is copied back to the initial signal node, so references don't have to be updated. This means that the memory management of Rizzo must be able to immediately identify signals that are no longer needed, otherwise these lead to unnecessary work being done. One solution to this is a reference counted memory management system. With reference counting, each object holds a counter with the number of references that exist for that object. The count then makes it possible to immediately know when a signal is not needed and can be safely discarded.

